\documentclass[11pt]{article}

\usepackage[utf8]{inputenc}
\usepackage[T1]{fontenc}
\usepackage[english]{babel}
\usepackage[a4paper,landscape,margin=1.55cm]{geometry}
\usepackage{booktabs}
\usepackage{longtable}
\usepackage{array}
\usepackage{enumitem}
\usepackage{hyperref}
\usepackage{xcolor}

\hypersetup{
  colorlinks=true,
  linkcolor=blue!55!black,
  urlcolor=blue!55!black
}

\newcommand{\coef}[2]{#1; $p=#2$}
\newcommand{\sig}{\textit{p}<0.001}

\title{Platform Membership, Marketing-Mix Effectiveness, and Capacity Utilisation\\
\large Model Tests, Extended Features, and Interpretation}
\author{JAMS Platform Membership Analysis}
\date{\today}

\begin{document}
\maketitle

\section{Purpose}

This document summarizes the current empirical evidence for the platform membership
paper. The objective is not to test whether platform membership mechanically raises
capacity utilisation. Instead, the analysis tests whether platform membership changes
the effectiveness of marketing-mix instruments, especially pricing, promotion, place,
product, and process-related features. The central dependent variable is transformed
capacity utilisation, \texttt{ur\_boxcox}. Trip volume, \texttt{log\_trip\_count\_day},
is used as a secondary outcome for selected robustness checks.

The empirical panel is a city--operator--day panel for 2024, covering March 11, 2024
through August 22, 2024. The main full-sample models use 14,514 observations, 97
city--operator units, and 30 cities. Platform membership is coded as substantive
bookable/payable integration and separated from weak exposure. Platform types distinguish
large aggregator membership, local MaaS membership, and multi-platform membership.

\section{Model Families}

The empirical strategy uses a family of fixed-effect regressions with cluster-robust
standard errors at the city--operator level. The two main fixed-effect strategies are:

\begin{itemize}[leftmargin=*]
  \item \textbf{Market FE:} city, date, and operator fixed effects.
  \item \textbf{Unit FE:} city--operator and date fixed effects.
\end{itemize}

The market-FE models identify differences across operators and cities after controlling
for date shocks and operator baselines. The unit-FE models are stricter: they identify
within-city--operator changes over time net of date fixed effects. Robust findings are
therefore those that survive both market-FE and unit-FE specifications or are repeatedly
confirmed across threshold, city-drop, outcome, and feature-extension tests.

The core model classes are:

\begin{enumerate}[leftmargin=*]
  \item \textbf{Platform main-effect models:}
  \[
  Y_{i c t} = \beta_1 Platform_{i c t} + X_{i c t}\Gamma + \alpha + \epsilon_{i c t}.
  \]
  \item \textbf{Marketing-mix moderation models:}
  \[
  Y_{i c t} = \beta_1 M_{i c t} + \beta_2 Platform_{i c t}
  + \beta_3 M_{i c t} \times Platform_{i c t} + X_{i c t}\Gamma + \alpha + \epsilon_{i c t}.
  \]
  \item \textbf{Platform-type threshold models:}
  \[
  Y_{i c t} = \sum_k \beta_k HighUnlockFee_{i c t} \times PlatformType^k_{i c t}
  + X_{i c t}\Gamma + \alpha + \epsilon_{i c t}.
  \]
  \item \textbf{Extended feature models:}
  \[
  Y_{i c t} = Feature_{i c t} \times PlatformType_{i c t}
  + X_{i c t}\Gamma + \alpha + \epsilon_{i c t}.
  \]
\end{enumerate}

Controls include own price per minute, unlock fee, relative fleet size, coverage,
exclusive coverage, promotion activity, and weather controls, depending on the model
family. Extended feature tests additionally use same-mode competition, same-platform
competition, active-vehicle shares, relative price gaps, commute/night usage shares,
and externally coded product/process/pass features.

\section{Feature Construction}

The extended feature layer adds three kinds of variables.

\paragraph{Competition and market structure.}
Same-city and same-mode competitor counts measure direct competitive pressure.
Concentration is captured through trip-share and active-vehicle-share Herfindahl-type
indices. Same-mode measures are preferred because scooter-to-scooter competition is
more meaningful than a pooled count of all shared mobility providers.

\paragraph{Relative price position.}
The most important new features compare each provider's price to the cheapest same-mode
competitor in the same city on the same day. The key variable is
\texttt{unlock\_fee\_gap\_to\_same\_mode\_min\_z}. It measures whether the provider has
a salient upfront access-price disadvantage relative to direct same-mode alternatives.
A parallel variable is constructed for price per minute.

\paragraph{Product, process, and pass affordances.}
The panel already contains group ride, vehicle-type breadth, subscription options,
and complementary services. Additional provider-level external coding captures
reservation availability, pass/free-unlock affordances, minute bundles, and mandatory
parking/geofence rules. These external flags are useful for theory development, but
because they vary mostly at the provider level, they are not yet cleanly estimable as
causal city-date effects.

\section{Results}

\setlength{\tabcolsep}{2.2pt}
\renewcommand{\arraystretch}{1.18}
\small
\begin{longtable}{p{0.075\linewidth}p{0.185\linewidth}p{0.10\linewidth}p{0.115\linewidth}p{0.115\linewidth}p{0.17\linewidth}p{0.085\linewidth}}
\caption{Relevant Findings, Robustness, and Full-Sample Estimates}\\
\toprule
\textbf{Finding} & \textbf{Substantive claim} & \textbf{Robustness} &
\textbf{Market FE} & \textbf{Unit FE} & \textbf{Observation basis} & \textbf{Primary source}\\
\midrule
\endfirsthead
\toprule
\textbf{Finding} & \textbf{Substantive claim} & \textbf{Robustness} &
\textbf{Market FE} & \textbf{Unit FE} & \textbf{Observation basis} & \textbf{Primary source}\\
\midrule
\endhead
\bottomrule
\endfoot

F1 & Platform membership does not create a stable generic utilisation lift. &
Not robust as main effect &
\coef{0.061}{0.008} &
\coef{-0.010}{0.205} &
Full sample, $N=14{,}514$. Positive only in market FE; unit FE not significant. &
\texttt{idea\_tests/05} \\

F2 & Unlock fees become more detrimental in platform contexts. &
Robust in market FE; weaker in strict unit FE &
\coef{-0.0715}{<0.001} &
\coef{-0.0236}{0.349} &
Full sample, $N=14{,}514$; negative in 9 significant concept-family models. &
\texttt{concept\_aligned} \\

F3 & High unlock fees are harmful in platform contexts. &
Robust &
\coef{-0.104}{<0.001} &
\coef{-0.040}{<0.001} &
Full sample, $N=14{,}514$; negative in 6/6 threshold models. &
\texttt{deep\_dive} \\

F3/F4 & High unlock fee $\times$ Local MaaS. &
Very robust &
\coef{-0.150}{0.002} &
\coef{-0.046}{<0.001} &
Full sample, $N=14{,}514$; negative in 6/6 models. &
\texttt{deep\_dive} \\

F3/F4 & High unlock fee $\times$ Multi-platform. &
Very robust &
\coef{-0.131}{0.007} &
\coef{-0.098}{<0.001} &
Full sample, $N=14{,}514$; negative in 6/6 models. &
\texttt{deep\_dive} \\

F3 & EUR 1.00 unlock-fee threshold $\times$ Local MaaS. &
Very robust &
\coef{-0.229}{<0.001} &
\coef{-0.049}{0.002} &
Full sample, $N=14{,}514$; Local MaaS high-fee support: 90 rows. &
\texttt{threshold} \\

F3 & EUR 1.00 unlock-fee threshold $\times$ Multi-platform. &
Very robust &
\coef{-0.146}{<0.001} &
\coef{-0.099}{<0.001} &
Full sample, $N=14{,}514$; multi-platform high-fee support: 1,308 rows. &
\texttt{threshold} \\

F3 & EUR 1.20 unlock-fee threshold $\times$ Local MaaS. &
Very robust &
\coef{-0.233}{<0.001} &
\coef{-0.047}{0.006} &
Full sample, $N=14{,}514$; Local MaaS high-fee support: 72 rows. &
\texttt{threshold} \\

F3 & EUR 1.20 unlock-fee threshold $\times$ Multi-platform. &
Very robust &
\coef{-0.171}{<0.001} &
\coef{-0.082}{<0.001} &
Full sample, $N=14{,}514$; multi-platform high-fee support: 1,029 rows. &
\texttt{threshold} \\

F3a & Local MaaS high-fee penalty is not driven by one platform city. &
Robust &
Median city-drop coefficient at EUR 1.00: -0.209 &
Median city-drop coefficient at EUR 1.00: -0.035 &
13/13 leave-one-platform-city models negative and significant for EUR 0.99, 1.00, and 1.20. &
\texttt{city\_sensitivity} \\

F3a & Multi-platform high-fee penalty is not a single-city artifact. &
Robust at EUR 1.00 and EUR 1.20 &
Median city-drop coefficients: -0.107 at EUR 1.00; -0.135 at EUR 1.20 &
Median city-drop coefficients: -0.085 at EUR 1.00; -0.073 at EUR 1.20 &
13/13 leave-one-platform-city models negative and significant at EUR 1.00 and 1.20. &
\texttt{city\_sensitivity} \\

F4a & Large aggregator-only / Free Now-style membership does not show a stable high-fee penalty. &
Negative evidence &
\coef{-0.045}{0.066} &
\coef{-0.016}{0.325} &
Full sample, $N=14{,}514$; not robustly significant. &
\texttt{threshold} \\

F3b & Local MaaS high-fee result generalizes from utilisation to trip volume. &
Robust for Local MaaS &
Log trips, EUR 1.00: \coef{-1.014}{<0.001} &
Log trips, EUR 1.00: \coef{-0.130}{<0.001} &
Full sample, $N=14{,}514$; not a pure utilisation-ratio artifact. &
\texttt{outcome} \\

F3b & Multi-platform trip-volume evidence is weaker than utilisation evidence. &
Moderate &
Log trips, EUR 1.00: \coef{-0.538}{<0.001} &
Log trips, EUR 1.00: \coef{-0.010}{0.873} &
Full sample, $N=14{,}514$; strong in market FE, not in unit FE. &
\texttt{outcome} \\

F5 & Generic promotion $\times$ platform membership is positive mainly within unit. &
Moderate &
\coef{0.012}{0.424} &
\coef{0.0115}{<0.001} &
Full sample, $N=14{,}514$; promotion rows: 738. &
\texttt{promotion} \\

F6 & Subscriber-free-minutes promotions fit multi-platform exposure. &
Exploratory mechanism, locally bounded &
\coef{0.122}{<0.001} &
\coef{0.043}{<0.001} &
Full sample model, $N=14{,}514$, but treated support is only Nürnberg--VOI, 144 rows. &
\texttt{deep\_dive} \\

F16 & Vehicle-type breadth buffers Local MaaS high-fee penalties. &
Moderate to robust &
\coef{0.074}{<0.001} &
\coef{0.014}{<0.001} &
Full sample, $N=14{,}514$; Local MaaS high-fee support: 90 rows. &
\texttt{feature\_fee} \\

F19 & Relative same-mode unlock-fee disadvantage $\times$ Local MaaS. &
Very robust; strongest new feature result &
\coef{-0.089}{<0.001} &
\coef{-0.030}{0.002} &
Full sample, $N=14{,}514$; feature fully observed. &
\texttt{extended} \\

F19 & Relative same-mode unlock-fee disadvantage $\times$ Multi-platform. &
Very robust &
\coef{-0.091}{<0.001} &
\coef{-0.041}{<0.001} &
Full sample, $N=14{,}514$; feature fully observed. &
\texttt{extended} \\

F20 & Relative per-minute price disadvantage is weaker than relative unlock-fee disadvantage. &
Moderate &
\coef{-0.104}{0.003} &
\coef{-0.020}{0.153} &
Full sample, $N=14{,}514$; weaker and less robust than unlock-fee gap. &
\texttt{extended} \\

F21 & Same-mode competition is a Local MaaS boundary condition, not the main mechanism. &
Moderate &
\coef{-0.059}{0.013} &
\coef{-0.019}{0.083} &
Full sample, $N=14{,}514$; same-mode competitor feature fully observed. &
\texttt{extended} \\

F22 & Own same-mode active-vehicle share helps platform contexts but does not offset high fees. &
Moderate &
For multi-platform: \coef{-0.007}{0.819} &
For multi-platform: \coef{0.019}{0.024} &
Full sample, $N=14{,}514$; positive mainly within unit. &
\texttt{extended} \\

F23 & Local MaaS performs worse in night-oriented usage contexts. &
Moderate &
\coef{-0.019}{0.001} &
\coef{-0.018}{<0.001} &
Full sample, $N=14{,}514$; night-trip share from raw trips. &
\texttt{extended} \\

F24 & Provider-level pass, free-unlock, reservation, and parking/geofence flags are useful but not robust estimators. &
Identification caveat &
Not interpreted due to rank/collinearity issues &
Not interpreted due to weak within-unit variation &
Full sample fields available, but scooter-only pass/minute-bundle support has only one unique value. &
\texttt{extended} \\

F25 & Same-platform competitor count is not estimable cleanly. &
Not supported &
Implausibly large coefficients and missing p-values &
Unstable and collinear &
Feature is mechanically tied to platform/city/operator structure. &
\texttt{extended} \\

\end{longtable}
\normalsize

\section{Neutral Description of Results}

The evidence rejects a simple main-effect interpretation of platform membership.
Platform membership is positive in market-FE models, but the same coefficient becomes
small and insignificant in unit-FE models. This pattern suggests that platform membership
is not best understood as a uniform utilisation booster. Instead, the stronger empirical
regularities appear in interaction terms between platform membership and marketing-mix
variables.

The most consistent baseline result is the pricing-friction effect. Unlock fees become
more detrimental in platform contexts, especially once the fee crosses salient high-fee
thresholds around EUR 1.00 and EUR 1.20. The threshold results are strongest for Local
MaaS and multi-platform contexts. At the EUR 1.00 cutoff, Local MaaS has a market-FE
coefficient of -0.229 and a unit-FE coefficient of -0.049. Multi-platform membership
has a market-FE coefficient of -0.146 and a unit-FE coefficient of -0.099. Both effects
are statistically significant. Similar estimates appear at EUR 1.20.

The leave-one-city tests show that these high-fee platform-type results are not driven
by a single city. For Local MaaS, all 13 leave-one-platform-city tests remain negative
and significant for EUR 0.99, EUR 1.00, and EUR 1.20. For multi-platform membership,
the EUR 1.00 and EUR 1.20 thresholds also remain negative and significant in all 13
city-drop models. The effect is therefore not only a Berlin, Hamburg, or Munich artifact.

Outcome robustness further strengthens the Local MaaS pricing-friction result. The
Local MaaS high-fee penalty appears not only in capacity utilisation but also in trip
volume. For log trips, the EUR 1.00 Local MaaS effect is -1.014 in market FE and -0.130
in unit FE. Multi-platform evidence is cleaner for capacity utilisation than for trip
volume, because the log-trip unit-FE coefficient is not significant.

The extended feature models sharpen the interpretation. The relative unlock-fee gap to
the cheapest same-mode competitor is the strongest new feature. In Local MaaS contexts,
the interaction is -0.089 in market FE and -0.030 in unit FE. In multi-platform contexts,
it is -0.091 in market FE and -0.041 in unit FE. These estimates indicate that platform
contexts are especially punitive when a provider has a higher upfront access fee than
direct same-mode alternatives.

Other features are weaker or more bounded. Same-mode competition is negative for Local
MaaS in market FE and marginal in unit FE, but it is less robust than relative unlock-fee
position. Per-minute price gaps are also weaker than unlock-fee gaps. Active-vehicle
share helps in some platform contexts, especially within unit, but it does not eliminate
high-fee penalties. Night-oriented usage contexts are negatively associated with Local
MaaS platform effectiveness. Promotion results are promising but locally bounded:
subscriber-free-minutes perform well in multi-platform models but are based on one
treated city--operator cell, Nürnberg--VOI.

\section{Interpretation and Analysis}

The results support a comparison-salience theory of platform membership. Platforms do
not simply create demand by adding another channel. Instead, they change the information
environment in which users evaluate mobility options. When multiple providers or modes
are visible in an integrated interface, upfront price differences become more salient.
This is why the unlock-fee effects are stronger and more consistent than per-minute
price effects. Unlock fees are immediate, discrete, and easy to compare before the trip
begins. Per-minute prices are less salient because their total cost implications are
uncertain until trip duration is known.

Local MaaS and multi-platform memberships are the key contexts. Large aggregator-only
membership does not show the same stable high-fee penalty. This matters theoretically:
the relevant mechanism is not simply commercial aggregator exposure. Rather, the strongest
effects occur when shared mobility is embedded in a broader urban mobility choice set.
Local MaaS interfaces and multi-platform exposure plausibly make outside options more
visible and comparable. In that setting, a high or relatively high unlock fee becomes a
stronger deterrent.

The new relative-price features substantially improve the story. Absolute high-fee
thresholds show that salient fees matter. Relative same-mode unlock-fee gaps show why
they matter in a platform context: users do not evaluate the fee in isolation; they can
compare it against direct alternatives. This provides a more marketing-specific account
of platform membership. The platform does not merely moderate price sensitivity; it
changes the competitive frame through which price is interpreted.

The vehicle-type breadth result adds a useful product-side counterpoint. Broader product
assortment appears to buffer Local MaaS high-fee penalties. This suggests that marketing
mix elements can interact: pricing friction is harmful, but product breadth can partially
offset it by improving modal fit. The finding should not be overstated as a universal
product-breadth effect, but it gives the paper a more balanced story than ``platforms
make fees bad.'' Platforms amplify pricing frictions, yet product design can soften the
effect.

Promotion is theoretically aligned with the mechanism but empirically less clean. The
subscriber-free-minutes result suggests that usage-cost-reducing promotions can work
well in multi-platform contexts. However, the evidence is localized to Nürnberg--VOI.
Therefore, promotion should be used as mechanism evidence rather than as a central causal
claim. Future data with campaign start and end dates across multiple cities would allow
a stronger test.

The competition findings are also better interpreted as boundary conditions than as the
main effect. Same-mode competition matters moderately in Local MaaS contexts, but the
relative unlock-fee gap is stronger. This implies that competition per se is not enough.
What matters is whether competition makes a specific marketing-mix disadvantage visible.
In other words, competition enables comparison, but the marketing-mix attribute being
compared remains central.

\section{Implications for the Paper}

The strongest paper-facing claim is:

\begin{quote}
Platform membership reshapes marketing-mix effectiveness by increasing the salience of
comparable mobility alternatives. In Local MaaS and multi-platform contexts, providers
with high or relatively high unlock fees experience lower capacity utilisation. This
effect is not a generic platform main effect, not a Free Now-only result, and not driven
by a single city.
\end{quote}

The main empirical table should therefore prioritize: (1) high unlock-fee thresholds
for Local MaaS and multi-platform membership; (2) relative same-mode unlock-fee gaps;
(3) outcome robustness for trip volume; and (4) city-drop robustness. Product breadth
can be used as a secondary moderation result. Promotion and provider-level pass features
should be framed as mechanism evidence and future data opportunities unless richer
city-date coding is added.

\section{Limitations}

Several limitations should be explicit. First, the data are observational and city--
operator--day based. The results identify robust conditional associations rather than
randomized treatment effects. Second, some platform types are structurally tied to cities
and providers, making unit-FE evidence especially important. Third, provider-level external
features such as pass availability, reservation, and parking/geofence rules lack enough
within-unit time variation for robust estimation. Fourth, promotion results are based on
few treated campaign contexts. Finally, same-platform competitor count is currently too
collinear with platform membership and city/operator structure to be interpreted as a
clean moderator.

\end{document}
